\documentclass[10pt]{article}
\usepackage[utf8]{inputenc}
\usepackage[T1]{fontenc}
\usepackage{graphicx}
\usepackage[export]{adjustbox}
\graphicspath{ {./images/} }
\usepackage{amsmath}
\usepackage{amsfonts}
\usepackage{amssymb}
\usepackage[version=4]{mhchem}
\usepackage{stmaryrd}
\usepackage{caption}

\begin{document}
\captionsetup{singlelinecheck=false}
Imperfect information

\section*{Refinements-Imperfect information}
\begin{itemize}
  \item Refinements of extensive form games
  \item Weakly Perfect Bayesian Equilibrium (WPBE):
  \item Sequential Equilibrium (SE)
  \item Forward induction concepts.
\end{itemize}

\section*{Non-credible threats, perfect information}
\begin{center}
\includegraphics[max width=\textwidth, alt={}]{fd69804c-f915-467d-b964-ffa616621981-04_515_907_444_190}
\end{center}

\begin{itemize}
  \item Two Nash equilibria
  \item $(d, T)$.
  \item $(u, B)$. Not SPE. Subgame perfection eliminates noncredible threats in games of perfect information.
\end{itemize}

\section*{Threats and imperfect information}
\section*{Example 1}
\begin{center}
\includegraphics[max width=\textwidth, alt={}]{fd69804c-f915-467d-b964-ffa616621981-05_1277_994_504_153}
\end{center}

NE: $[(L, L), B], \quad\left(E U_{1}, E U_{2}\right)=(2,2)$\\
NE: $[(R, L), T], \quad\left(E U_{1}, E U_{2}\right)=(2.5,1.5)$\\
$[(L, L), B]$ is an unintuitive NE

\begin{itemize}
  \item This is a signaling game:
  \item Nature flips a coin.
  \item Player 1 observes either $h_{1}$ (head) or $h_{1}^{\prime}$ (tail), and chooses an action.
  \item Player 2 does not observe nature's move but observes player 1's action and responds.
  \item The game has no subgames, so every NE is SPE.
  \item Equilibrium $[(L, L, B)]$ is not intuitive.
  \item For player $2, B$ is a dominated strategy, and therefore it should not be credible. Subgame perfection does not apply.
\end{itemize}

WPBE

\section*{Assessment}
Fix a finite extensive-form game $\Gamma$.

\section*{Definition 1 (Assessment)}
\begin{itemize}
  \item A system of beliefs is a profile $\mu=\left(\mu_{1}, \ldots, \mu_{I}\right)$ such that
\end{itemize}

$$
(\forall i \in N \backslash\{0\})\left(\forall h \in \mathcal{U}^{i}\right) \quad \mu_{i}(\cdot \mid h) \in \Delta(h),
$$

where $\Delta(h)$ is the set of probability distributions on the information set $h$.

For $x \in h, \mu_{i}(x \mid h)$ is the probability assigned to $x$ by beliefs $\mu_{i}(\cdot \mid h)$.\\
An assessment is a strategy-profile and system-of-beliefs pair $(\sigma, \mu)$.

\section*{Weakly Perfect Bayesian Equilibrium}
Definition 2 (WPBE) An assessment $(\sigma, \mu)$ is a Weakly Perfect Bayesian Equilibrium (WPBE) of a finite, extensive-form game if

\begin{enumerate}
  \item Sequential rationality: $(\forall i \in N)\left(\forall h \in \mathcal{U}_{i}\right)\left(\forall \sigma_{i}^{\prime} \in \Sigma_{i}\right)$
\end{enumerate}

$$
E\left[u_{i} \mid h, \mu_{i}, \sigma_{i}, \sigma_{-i}\right] \geq E\left[u_{i} \mid h, \mu_{i}, \sigma_{i}^{\prime}, \sigma_{-i}\right]
$$

\begin{enumerate}
  \setcounter{enumi}{1}
  \item Bayes rule: $(\forall i \in N)\left(\forall h \in \mathcal{U}_{i}\right)$,
\end{enumerate}

$$
P_{\sigma}(h)>0 \Longrightarrow \forall x \in h, \mu_{i}(x \mid h)=\frac{P_{\sigma}(x)}{P_{\sigma}(h)}
$$

\section*{NE and sequential rationality}
Fix a finite, extensive-form game $\Gamma$ of perfect recall.\\
Theorem 1 Let $(\sigma, \mu)$ be an assessment. If at every information set in the path of play,\\
i. $\mu$ is derived by Bayes rule from $\sigma$ and\\
ii. $(\sigma, \mu)$ satisfies sequential rationality, then $\sigma$ is a NE of $\Gamma$.

Theorem 2 Let $\sigma$ be a NE and $\mu$ a system of beliefs derived by Bayes rule from $\sigma$ when possible. Then $(\sigma, \mu)$ is sequentially rational at every information set on the equilibrium path.

\begin{itemize}
  \item Maschler, Solan, and Zamir (2013), Theorems 7.44, 7.45. pp. 276-277.
\end{itemize}

\section*{Sequential rationality. Interpretation}
\begin{itemize}
  \item Subgame perfection eliminates non-credible threats by requiring that a NE be played in any subgame.
  \item Games of imperfect information may have no subgames. A ˋ"game" starting at any information set may be defined.
  \item Interpret an information set $h$ as the beginning of the game; $\mu_{i}(\cdot \mid h)$ as the probability of being at each $x \in h$.
  \item Available actions remain as in the original game.
  \item $E\left[u_{i} \mid h, \mu_{i}, \sigma_{i}, \sigma_{-i}\right]$ denotes $i$ 's expected payoff in the "game" that begins at information set $h$.
\end{itemize}

\section*{Example to illustrate WPBE 1}
\begin{center}
\includegraphics[max width=\textwidth, alt={}]{fd69804c-f915-467d-b964-ffa616621981-12_1390_2179_348_150}
\end{center}

\section*{Player 1's strategy and beliefs 2}
$$
\begin{aligned}
\sigma_{1}\left(h_{1}\right) & =\underbrace{(0.5,0.5)}_{(U, D)} \\
\sigma_{1}\left(h_{1}^{\prime}\right) & =\overbrace{(1,0)} \\
\mu_{1}\left(h e a d \mid h_{1}\right) & =1 \\
\mu_{1}\left(t a i l \mid h_{1}^{\prime}\right) & =1
\end{aligned}
$$

\section*{Player 2's strategy and beliefs 3}
$$
\begin{aligned}
& \sigma_{2}\left(h_{2}\right)=\underbrace{(0.5,0.5)}_{(R, L)} \\
& \sigma_{2}\left(h_{2}^{\prime}\right)=\overbrace{(1,0)}
\end{aligned}
$$

Consider the beliefs:

$$
\begin{aligned}
\mu_{2}\left((\text { head }, U) \mid h_{2}\right) & =0.5 \\
\mu_{2}\left((\text { tail }, U) \mid h_{2}^{\prime}\right) & =1
\end{aligned}
$$

\section*{Revisit Example 1}
\begin{center}
\includegraphics[max width=\textwidth, alt={}]{fd69804c-f915-467d-b964-ffa616621981-15_1281_1007_391_149}
\end{center}

NE: $[(R, L), T], \quad\left(E U_{1}, E U_{2}\right)=(2.5,1.5)$\\
1's beliefs: $\mu_{1}\left(\right.$ head $\left.\mid h_{1}\right)=1, \mu_{1}\left(\right.$ tail $\left.\left.\mid h_{1}^{\prime}\right)\right)=1$\\
2's beliefs: $\mu_{2}\left((\right.$ head,$\left.R) \mid h_{2}\right)=1$

\begin{itemize}
  \item 2's beliefs: $h_{2}=\{($ head,$R),($ tail,$R)\}$.
\end{itemize}

$$
\begin{aligned}
& \mu_{2}\left((\text { head }, R) \mid h_{2}\right) \\
& \quad=\frac{P_{\sigma}((\text { head }, R))}{P_{\sigma}\left(h_{2}\right)} \\
& \quad=\frac{P_{0}(\text { head }) \times \sigma_{1}\left(R \mid h_{1}\right)}{P_{0}(\text { head }) \times \sigma_{1}\left(R \mid h_{1}\right)+P_{0}(\text { tail }) \times \sigma_{1}\left(R \mid h_{1}^{\prime}\right)} \\
& \quad=\frac{1 / 2 \times 1}{1 / 2 \times 1+1 / 2 \times 0}=\frac{1 / 2}{1 / 2}=1
\end{aligned}
$$

\begin{itemize}
  \item 2's sequential rationality follows. See previous figure.
  \item 2's sequential rationality if not yet convinced.\\
$E\left[U_{2} \mid h_{2}, T, \sigma_{1}\right]=\mu_{2}\left((\right.$ head,$\left.R) \mid h_{2}\right) 1+\mu_{2}\left((\right.$ tail,$\left.R) \mid h_{2}\right) 0=1$\\
$E\left[U_{2} \mid h_{2}, B, \sigma_{1}\right]=\mu_{2}\left((\right.$ head,$\left.R) \mid h_{2}\right) 0+\mu_{2}\left((\right.$ tail,$\left.R) \mid h_{2}\right)(-1)=0$\\
$E\left[U_{2} \mid h_{2}, T, \sigma_{1}\right]=1>0=E\left[U_{2} \mid h_{2}, B, \sigma_{1}\right]$
\end{itemize}

\section*{The other equilibrium in Example 1}
\begin{center}
\includegraphics[max width=\textwidth, alt={}]{fd69804c-f915-467d-b964-ffa616621981-18_1277_1002_391_149}
\end{center}

NE: $[(L, L), B], \quad\left(E U_{1}, E U_{2}\right)=(2,2)$.\\
$[(L, L), B]$ is the unintuitive NE.\\
$P_{\sigma}\left(h_{2}\right)=0$; then free to choose beliefs.\\
$\left(\forall \mu_{2}\left(x \mid h_{2}\right) \in[0,1]\right)$ 2's best action is T.\\
No beliefs support 2's choosing $B$ at $h_{2}$.\\
Hence $[(L, L), B]$ is not a WPBE.

\section*{Example: WPBE need not be SPE}
\begin{itemize}
  \item A potential entrant (E), decides whether to enter a market, (in) or (out).
  \item If E chooses in, then E and the incumbent (I), decide simultaneous whether to fight or acquiesce. Incumbent actions are in lower case.
\end{itemize}

\begin{figure}[h]
\begin{center}
  \includegraphics[alt={},max width=\textwidth]{fd69804c-f915-467d-b964-ffa616621981-19_577_1224_855_261}
\captionsetup{labelformat=empty}
\caption{Figure 1}
\end{center}
\end{figure}

\begin{itemize}
  \item Pure-strategy NE:
  \item ((OUT, F), f),
  \item ((OUT, A), f),
  \item $((I N, A), a)$.
  \item Subgame Perfect: $((I N, A), a)$.
  \item A peculiar WPBE. $((O U T, A), f) ; \mu_{2}\left((I N, F) \mid h_{I}\right)=1$.
  \item Normal form.
\end{itemize}

\begin{center}
\begin{tabular}{lcc}
 & $\mathbf{a}$ & $\mathbf{f}$ \\
\hline
OUT-A & 0,2 & 0,2 \\
\hline
OUT-F & 0,2 & 0,2 \\
\hline
IN-A & 3,1 & $-2,-2$ \\
\hline
IN-F & $1,-2$ & $-3,-1$ \\
\hline
\end{tabular}
\end{center}

\begin{itemize}
  \item Reduced Normal form.
\end{itemize}

\begin{center}
\begin{tabular}{lcc}
 & $\mathbf{a}$ & $\mathbf{f}$ \\
\hline
OUT- & 0,2 & 0,2 \\
\hline
IN-A & 3,1 & $-2,-2$ \\
\hline
\end{tabular}
\end{center}

\begin{itemize}
  \item The Entrant has two information sets, say $h_{E}$ and $h_{E}^{\prime}$, both singletons. The Incumbent has a single information set, $h_{I}$.
  \item WPBE (OUT-A,f)
\end{itemize}

$$
\begin{aligned}
\sigma_{E}\left(h_{E}\right) & =(1,0) ;(\mathrm{OUT}, \mathrm{IN}) \\
\sigma_{E}\left(h_{E}^{\prime}\right) & =(1,0) ;(\mathrm{A}, \mathrm{~F}) \\
\sigma_{I}\left(h_{I}\right) & =(0,1) ;(\mathrm{a}, \mathrm{f}) \\
\mu_{I}\left((I N, A) \mid h_{I}\right) & =0
\end{aligned}
$$

This is WPBE but is not SPE; strategies do not specify an equilibrium in the subgme.

\section*{Example: unintuitive WPBE}
\begin{figure}[h]
\begin{center}
  \includegraphics[alt={},max width=\textwidth]{fd69804c-f915-467d-b964-ffa616621981-23_1298_1122_344_203}
\captionsetup{labelformat=empty}
\caption{Figure 2}
\end{center}
\end{figure}

\begin{figure}[h]
\begin{center}
  \includegraphics[alt={},max width=\textwidth]{fd69804c-f915-467d-b964-ffa616621981-24_1304_1116_178_209}
\captionsetup{labelformat=empty}
\caption{Figure 3}
\end{center}
\end{figure}

\begin{figure}[h]
\begin{center}
  \includegraphics[alt={},max width=\textwidth]{fd69804c-f915-467d-b964-ffa616621981-25_1276_1116_271_209}
\captionsetup{labelformat=empty}
\caption{Figure 4}
\end{center}
\end{figure}

First WPBE (blue)

$$
\begin{array}{ll}
\sigma_{1}\left(h_{1}\right)=(0,1) ;(L, R) . & \mu_{1}\left(h d \mid h_{1}\right)=\frac{1}{2} \\
\sigma_{2}\left(h_{2}\right)=(0,1) ;(T, B) . & \mu_{2}\left(h d, R \mid h_{2}\right)=\frac{1}{2}
\end{array}
$$

Second WPBE (red)

$$
\begin{array}{ll}
\sigma_{1}\left(h_{1}\right)=(1,0) ;(L, R) . & \mu_{1}\left(h d \mid h_{1}\right)=\frac{1}{2} \\
\sigma_{2}\left(h_{2}\right)=(1,0) ;(T, B) . & \mu_{2}\left(h d, R \mid h_{2}\right)=\frac{9}{10}
\end{array}
$$

\section*{Remarks}
\begin{itemize}
  \item Many WPBE because it imposes no restrictions on beliefs off the equilibrium path.
  \item Including beliefs in the definition of equilibrium possibilitates further refinements.
  \item Red equilibrium and the common prior assumption.
  \item Finding intuitive properties with examples, define solution concept. Weak as method.
\end{itemize}

\section*{References}
Maschler, Michael, Eilon Solan, and Shmuel Zamir. 2013. Game Theory. Translated by Ziv Hellman. 1st ed. Cambridge University Press.


\end{document}